\begin{textsample}{POS Dim 1 – vem\_ed\_33 – Score 19.00 – vem\_ed\_33/t180.txt}  \label{ex:f1_pos_077}
Receitas sem fronteiras: um PCI entre Brasil e Colômbia Artigo de \textbf{Opinião} A globalização mostra \textbf{que} o ensino \textbf{superior} tecnológico precisa ir além das fronteiras.
No Centro Paula Souza, a prática de Internacionalização em \textbf{Casa} \textbf{é} uma \textbf{realidade} graças aos Projetos Colaborativos Internacionais (PCIs) realizados sob orientação da Área de Internacionalização do Ensino Superior (AIES) da Divisão de Extensão e Pesquisa no Ensino Superior (Depes) da Coordenadoria Geral de Ensino Superior (CGESG).
Os PCIs \textbf{seguem} a abordagem internacionalmente conhecida como COIL (Collaborative Online International Learning, ou Aprendizagem Internacional Colaborativa Online).
Esses projetos ajudam no \textbf{desenvolvimento} de \textbf{competências} globais sem \textbf{que} \textbf{seja} necessário viajar.
Isso ocorre porque a colaboração internacional aumenta a \textbf{empatia} e a \textbf{capacidade} de se adaptar a novas situações, o \textbf{que} prepara o aluno para trabalhar em \textbf{ambientes} multiculturais mediados por novas tecnologias.
Este artigo apresenta uma experiência de Internacionalização em Casa em um Projeto Colaborativo Internacional (PCI), \textbf{que} ocorreu no segundo semestre de 2025 entre a Fatec São Roque e a Uniminuto (Colômbia).
Intitulado “Receitas sem fronteiras: videorecetario colombobrasileño”, o PCI envolveu meus alunos de espanhol, matriculados no curso de Gestão de Empreendimentos Gastronômicos da Fatec São Roque, e estudantes de Tecnologia em Realização Audiovisual da Uniminuto orientados pelas professoras Sandra Sanchez e Lorena Parodi.
Os fatecanos fizeram uma pesquisa sobre receitas típicas brasileiras e o desafio pedagógico consistiu em traduzir técnicas dos preparos e enviar instruções em áudio em espanhol para os colegas colombianos.
Tais instruções deveriam \textbf{ser} claras para \textbf{que} os estudantes da Uniminuto pudessem \textbf{seguir} e preparar as receitas, registrando-as em um vídeo curto.
As reuniões \textbf{síncronas} de quebra-gelo e apresentações finais das vídeoreceitas ocorreram no Teams da Uniminuto.
Ao \textbf{longo} da realização do projeto, os alunos se comunicavam em grupos de WhatsApp para discutir o \textbf{desenvolvimento} dos trabalhos.
A \textbf{produção} conjunta \textbf{foi} registrada e \textbf{publicada} na plataforma Padlet, \textbf{que} funciona como um mural digital \textbf{colaborativo}, apto para receber diversos tipos de arquivos (texto, imagem, links para vídeos).
O Padlet \textbf{representou} grande utilidade, pois \textbf{permitiu} o compartilhamento dos arquivos e os comentários às \textbf{produções} dos grupos mistos de \textbf{forma} rápida e simples.
A gastronomia brasileira \textbf{foi} a \textbf{ponte} necessária para a comunicação em espanhol (o nosso principal desafio) e a prova de \textbf{que} a cooperação técnica supera dificuldades.
A culinária brasileira teve um \textbf{impacto} muito forte nos colombianos.
Mesmo à distância, a comida \textbf{é} uma ferramenta poderosa para unir \textbf{pessoas} de diferentes países e culturas.
Uma prova de \textbf{que} essa receita de colaboração internacional deu certo \textbf{é} \textbf{que} em 2026 partimos para a segunda edição do projeto: um prato cheio de novas descobertas.

% matched lemmas: ambiente, capacidade, casa, colaborativo, competência, desenvolvimento, empatia, forma, impacto, longo, opinião, permitir, pessoa, ponte, produção, publicar, que, realidade, representar, seguir, ser, superior, síncrono
\end{textsample}
